\chapter{Pengintegralan pada Sebuah Ruas}\label{ch:b1:integration}

\omterm{thm:b1:integration:def}{Integral} pada jilid Sekolah Menengah dahulu didasarkan pada luas yang diambil
secara intuitif. Bab ini mengonstruksinya: mula-mula bagi \omterm{def:b1:integration:step}{fungsi tangga},
yang di situ \omterm{thm:b1:integration:def}{integralnya} berupa jumlah yang hingga, lalu bagi fungsi yang \omterm{def:b1:continuity:continuous}{kontinu} (dan
\omterm{def:b1:continuity:continuous}{kontinu} sepotong-sepotong) lewat hampiran seragam --- yaitu
tempat teorema Heine (\cref{thm:b1:continuity:heine}) mendapatkan
nafkahnya. Lalu teorema fundamental kalkulus menghubungkan
konstruksinya dengan primitif, sedangkan \omterm{thm:b1:integration:riemann}{jumlah Riemann} menghubungkannya dengan rata-rata
yang diskret.

Di sepanjang bab ini, $a < b$ merupakan bilangan real.

\section{Fungsi tangga}

\begin{definition}\label{def:b1:integration:step}
Fungsi $\varphi \colon \intcc{a}{b} \to \R$ disebut \emph{fungsi
tangga}\index{fungsi tangga} bila ada \omterm{thm:b1:logic:partition}{partisi} $a = x_0
< x_1 < \dots < x_n = b$ sedemikian sehingga $\varphi$ konstan, yang sama dengan
$c_i$, pada setiap \omterm{prop:b1:reals:intervals}{selang} \omterm{def:b1:topology:open}{terbuka} $\intoo{x_{i-1}}{x_i}$ (adapun nilainya di
simpulnya tak dibatasi). Adapun \omterm{thm:b1:integration:def}{integralnya} adalah
\[
\int_a^b \varphi = \sum_{i=1}^{n} c_i\,(x_i - x_{i-1}),
\]
yang tak bergantung pada \omterm{thm:b1:logic:partition}{partisi} yang dipilih (haluskanlah dua \omterm{thm:b1:logic:partition}{partisi} lewat
\omterm{thm:b1:logic:partition}{partisi} bersamanya: karena kedua ruasnya tak berubah oleh penghalusan).
\end{definition}

\begin{proposition}\label{prop:b1:integration:stepprops}
Pada \omterm{def:b1:integration:step}{fungsi tangga}, \omterm{thm:b1:integration:def}{integralnya} bersifat linear, naik (karena $\varphi \leq
\psi \implies \int\varphi \leq \int\psi$), dan memenuhi relasi
Chasles $\int_a^b = \int_a^c + \int_c^b$ untuk $a < c < b$.
\end{proposition}

\begin{proof}
Mesinnya adalah \emph{kekekalan terhadap penghalusan}, yang dinyatakan pada
definisinya: bahwa menyisipkan satu simpul tambahan $t \in \intoo{x_{i-1}}{x_i}$
ke dalam sebuah \omterm{thm:b1:logic:partition}{partisi} mengganti suku $c_i(x_i - x_{i-1})$ dengan
$c_i(t - x_{i-1}) + c_i(x_i - t)$ --- yaitu bilangan yang sama --- sehingga
\omterm{thm:b1:integration:def}{integralnya} tak berubah oleh penghalusan hingga mana pun. Sekarang ambillah
$\varphi$ dengan \omterm{thm:b1:logic:partition}{partisi} $\sigma$ dan $\psi$ dengan \omterm{thm:b1:logic:partition}{partisi}
$\sigma'$: maka pada penghalusan bersamanya $\sigma \cup \sigma'$ keduanya
merupakan \omterm{def:b1:integration:step}{fungsi tangga} dengan simpul yang \emph{sama}, dan pada setiap kepingannya
$\varphi + \lambda\psi$ konstan sama dengan $c_i + \lambda
d_i$: sehingga kelinearannya tereduksi menjadi kelinearan jumlah yang hingga. Adapun kenaikannya:
$c_i \leq d_i$ pada setiap kepingannya memberikan $\sum c_i \Delta_i \leq \sum
d_i \Delta_i$ (karena panjangnya $\Delta_i \geq 0$). Sedangkan Chasles: sisipkanlah
simpul $c$ lalu pecahlah jumlahnya di situ.
\end{proof}

\section{Integral fungsi yang kontinu}

\begin{theorem}[Hampiran seragam]\label{thm:b1:integration:approx}
Misalkan $f$ \omterm{def:b1:continuity:continuous}{kontinu} pada $\intcc{a}{b}$. Untuk setiap $\varepsilon > 0$
ada \omterm{def:b1:integration:step}{fungsi tangga} $\varphi, \psi$ dengan
\[
\varphi \leq f \leq \psi
\qquad\text{dan}\qquad
\psi - \varphi \leq \varepsilon \text{ pada } \intcc{a}{b}.
\]
\end{theorem}

\begin{proof}
Menurut teorema Heine (\cref{thm:b1:continuity:heine}), $f$ bersifat \omterm{def:b1:continuity:uniform}{kontinu
seragam}: pilihlah $\delta$ bagi $\varepsilon$, dan sebuah \omterm{thm:b1:logic:partition}{partisi}
berjala $< \delta$ (misalnya yang berjarak sama, dengan $n > \frac{b -
a}{\delta}$ keping). Pada setiap keping \omterm{def:b1:topology:closed}{tertutupnya} $\intcc{x_{i-1}}{x_i}$, $f$
mencapai minimum $m_i$ dan maksimum $M_i$
(\cref{thm:b1:continuity:evt}), dan $M_i - m_i \leq \varepsilon$
(karena kedua titik ekstremnya berjarak $\delta$). Definisikanlah $\varphi = m_i$
dan $\psi = M_i$ pada $\intoo{x_{i-1}}{x_i}$ (serta $\varphi = \psi = f$
di simpulnya).
\end{proof}

\begin{theorem}[Definisi integralnya]\label{thm:b1:integration:def}
Misalkan $f$ \omterm{def:b1:continuity:continuous}{kontinu} pada $\intcc{a}{b}$. Kedua bilangan
\[
I_-(f) = \sup\Bigl\{\int_a^b \varphi : \varphi \text{ tangga},\
\varphi \leq f\Bigr\},
\qquad
I_+(f) = \inf\Bigl\{\int_a^b \psi : \psi \text{ tangga},\ \psi \geq
f\Bigr\}
\]
bernilai sama; dan nilai bersamanya adalah
\emph{integral}\index{integral} $\int_a^b f$ (yang juga ditulis $\int_a^b
f(t)\,\dd t$). Ia berimpit dengan gagasan sebelumnya pada \omterm{def:b1:integration:step}{fungsi
tangga}, dan meluas ke fungsi yang \omterm{def:b1:continuity:continuous}{kontinu} sepotong-sepotong lewat
pemecahan $\intcc{a}{b}$ di titik ketakkontinuannya (dengan Chasles sebagai
definisinya di sana).
\end{theorem}

\begin{proof}
Kedua \omterm{def:b1:logic:sets}{himpunannya} tak kosong (karena $f$ terbatas) dan setiap \omterm{thm:b1:integration:def}{integral} tangga bawahnya
bernilai $\leq$ setiap yang atasnya (menurut kemonotonan pada \omterm{def:b1:integration:step}{fungsi tangga}): sehingga
$I_-(f) \leq I_+(f)$. Lalu menurut \cref{thm:b1:integration:approx}, untuk setiap
$\varepsilon$ ada sepasang dengan $\int\psi - \int\varphi \leq
\varepsilon(b - a)$: sehingga sup dan infnya terapit bersama, jadi $I_- =
I_+$.
\end{proof}

\begin{example}[Kontinu sepotong-sepotong, tanpa drama]\label{ex:b1:integration:piecewise}
\omterm{thm:b1:reals:floor}{Fungsi lantai} pada $\intcc{0}{3}$ merupakan \omterm{def:b1:integration:step}{fungsi tangga} yang
menyamar: karena dengan memecahnya di lompatannya,
\[
\int_0^3 \lfloor t \rfloor\,\dd t
= \int_0^1 0 + \int_1^2 1 + \int_2^3 2 = 0 + 1 + 2 = 3 ,
\]
dan nilainya \emph{di} titik lompatannya $1, 2$ tak relevan:
karena mengubah sebuah fungsi di titik yang berhingga banyak tak mengubah \omterm{thm:b1:integration:def}{integralnya}
(sebab \omterm{def:b1:integration:step}{fungsi tangga} pengapitnya tak terpengaruh). Inilah seluruh
isi perluasan ``\omterm{def:b1:continuity:continuous}{kontinu} sepotong-sepotong'' itu: potonglah di
titik ketakkontinuannya yang berhingga banyak, integralkanlah setiap keping yang \omterm{def:b1:continuity:continuous}{kontinu},
lalu jumlahkan --- jadi Chasles sebagai definisi.
\end{example}

\begin{example}[Definisinya menghitung, sekali]\label{ex:b1:integration:defcompute}
Misalkan $f(x) = x$ pada $\intcc{0}{1}$ lalu potonglah menjadi $n$ keping yang sama.
Adapun \omterm{def:b1:integration:step}{fungsi tangga} terbaik yang konstan pada kepingnya adalah $\varphi =
\frac{k-1}{n}$ dan $\psi = \frac kn$ pada keping ke-$k$, dengan
\[
\int_0^1 \varphi = \sum_{k=1}^{n} \frac{k-1}{n}\cdot\frac1n
= \frac{n-1}{2n},
\qquad
\int_0^1 \psi = \sum_{k=1}^{n} \frac{k}{n}\cdot\frac1n
= \frac{n+1}{2n} .
\]
Setiap \omterm{thm:b1:integration:def}{integral} bawahnya $\leq I_-(f) \leq I_+(f) \leq$ setiap
yang atasnya, sehingga $\frac{n-1}{2n} \leq I_-(f) \leq I_+(f) \leq
\frac{n+1}{2n}$ untuk setiap $n$: jadi keduanya terapit menuju $\frac12$, dan
$\int_0^1 x\,\dd x = \frac12$ langsung dari definisinya. Inti
gagasan penutupnya: bahwa inilah pertama dan terakhir kalinya kita mengintegralkan
dari definisinya --- karena teorema fundamental di bawah mengganti
semua perhitungan semacam itu dengan satu pencarian antiturunan, dan itulah
seluruh inti ekonomi bab ini.
\end{example}

\begin{theorem}[Sifatnya]\label{thm:b1:integration:props}
Untuk $f, g$ yang \omterm{def:b1:continuity:continuous}{kontinu} (atau \omterm{def:b1:continuity:continuous}{kontinu} sepotong-sepotong) pada $\intcc{a}{b}$ dan
$\lambda \in \R$:
\begin{enumerate}
  \item kelinearan: $\int (f + \lambda g) = \int f + \lambda \int g$;
  \item kemonotonan: $f \leq g \implies \int_a^b f \leq \int_a^b g$;
        dan $\bigl|\int_a^b f\bigr| \leq \int_a^b \abs f \leq (b -
        a)\, \sup\abs f$;
  \item Chasles: $\int_a^b = \int_a^c + \int_c^b$ (dengan
        kesepakatan $\int_b^a = -\int_a^b$, yang sah untuk urutan
        batas mana pun);
  \item kepositifan tegas: jika $f$ \omterm{def:b1:continuity:continuous}{kontinu}, $f \geq 0$ dan
        $\int_a^b f = 0$, maka $f = 0$ di mana-mana pada
        $\intcc{a}{b}$.
\end{enumerate}
\end{theorem}

\begin{proof}
(1)--(3) berpindah dari \omterm{def:b1:integration:step}{fungsi tangga} ke limitnya lewat
definisi sup/infnya. Adapun kelinearannya pantas diperinci sekali: diberikan
$\varepsilon > 0$, apitlah $\varphi_f \leq f \leq \psi_f$ dan
$\varphi_g \leq g \leq \psi_g$ dengan celah $\leq \varepsilon$
(\cref{thm:b1:integration:approx}). Untuk $\lambda \geq 0$,
$\varphi_f + \lambda\varphi_g \leq f + \lambda g \leq \psi_f +
\lambda\psi_g$ merupakan apitan oleh \omterm{def:b1:integration:step}{fungsi tangga} dengan celah $\leq
(1 + \lambda)\varepsilon$, dan \omterm{thm:b1:integration:def}{integral} tangganya sama dengan $\int
\varphi_f + \lambda\int\varphi_g$ dan seterusnya
(\cref{prop:b1:integration:stepprops}): sehingga membiarkan $\varepsilon \to
0$ mengapit $\int(f + \lambda g)$ menuju $\int f + \lambda\int g$.
Untuk $\lambda < 0$, mengalikan dengan $\lambda$ \emph{membalikkan}
apitan $g$ --- karena \omterm{def:b1:integration:step}{fungsi tangga} bawah bagi $\lambda g$ adalah
$\lambda \psi_g$ --- dan apitan yang sama berjalan dengan perannya
tertukar. Adapun batas $\abs{\int f} \leq \int \abs f$ datang dari
$-\abs f \leq f \leq \abs f$ dan kemonotonannya.

(4) Lewat kontraposisinya: jika $f(x_0) = m > 0$, maka \omterm{def:b1:continuity:continuous}{kekontinuannya} menyediakan
subselang berpanjang $\eta > 0$ yang di situ $f \geq \frac m2$; lalu
\omterm{def:b1:integration:step}{fungsi tangga} yang bernilai $\frac m2$ di sana dan $0$ di tempat lain bersifat $\leq f$,
sehingga $\int f \geq \frac{m\eta}{2} > 0$.
\end{proof}

\begin{example}[Chasles dalam kerja: integral dengan nilai mutlak]\label{ex:b1:integration:chasleswork}
Untuk mengintegralkan nilai mutlak, potonglah di tempat tandanya berubah.
\[
\int_0^2 \abs{x - 1}\,\dd x
= \int_0^1 (1 - x)\,\dd x + \int_1^2 (x - 1)\,\dd x
= \frac12 + \frac12 = 1 ,
\]
lalu, dengan memotong $\intcc{0}{2\pi}$ di $\pi$:
\[
\int_0^{2\pi} \abs{\sin t}\,\dd t
= \int_0^{\pi} \sin t\,\dd t - \int_{\pi}^{2\pi} \sin t\,\dd t
= 2 + 2 = 4 ,
\]
sedangkan $\int_0^{2\pi} \sin t\,\dd t = 0$: jadi peniadaannya nyata,
dan itulah sebabnya \omterm{def:b1:logic:statement}{pernyataan} kepositifan tegasnya
(\cref{thm:b1:integration:props}~(4)) membawa hipotesis $f
\geq 0$ --- karena tanpanya, \omterm{thm:b1:integration:def}{integral} yang lenyap tak membuktikan apa pun
tentang $f$. Inti gagasan penutupnya: bahwa $\int \abs f$ mengukur
\emph{luas}, sedangkan $\int f$ mengukur \emph{neraca bertanda}; sehingga
ketaksamaan $\abs{\int f} \leq \int\abs f$ merupakan catatan yang persis tentang
apa yang dapat dihancurkan peniadaannya.
\end{example}

\section{Teorema fundamental kalkulus}

\begin{theorem}[Teorema fundamental kalkulus]\label{thm:b1:integration:ftc}
Misalkan $f$ \omterm{def:b1:continuity:continuous}{kontinu} pada sebuah \omterm{prop:b1:reals:intervals}{selang} $I$ dan $a \in I$. Maka fungsi
\[
F(x) = \int_a^x f(t)\, \dd t
\]
berkelas $C^1$ pada $I$, dengan $F' = f$: jadi setiap fungsi yang \omterm{def:b1:continuity:continuous}{kontinu}
pada sebuah \omterm{prop:b1:reals:intervals}{selang} mempunyai primitif. Akibatnya, untuk \emph{sebarang}
primitif $G$ bagi $f$:\index{teorema fundamental kalkulus}
\[
\int_a^b f(t)\,\dd t = G(b) - G(a) .
\]
\end{theorem}

\begin{proof}
Tetapkanlah $x_0 \in I$ dan $\varepsilon > 0$; maka \omterm{def:b1:continuity:continuous}{kekontinuan} di $x_0$ menyediakan
$\delta$ dengan $\abs{f(t) - f(x_0)} \leq \varepsilon$ untuk $\abs{t -
x_0} \leq \delta$. Untuk $0 < \abs{h} \leq \delta$ (dan $x_0 + h \in
I$), Chasles memberikan
\[
\frac{F(x_0 + h) - F(x_0)}{h} - f(x_0)
= \frac 1h \int_{x_0}^{x_0+h} \bigl(f(t) - f(x_0)\bigr)\dd t ,
\]
yang nilai mutlaknya $\leq \frac{1}{\abs h}\cdot \abs h\,
\varepsilon = \varepsilon$ (menurut batas (2), yang sah untuk urutan
batas mana pun). Jadi $F'(x_0) = f(x_0)$; dan $F' = f$ \omterm{def:b1:continuity:continuous}{kontinu}: sehingga $F$
berkelas $C^1$. Jika $G' = f$ juga, maka $(G - F)' = 0$ pada \omterm{prop:b1:reals:intervals}{selangnya}, sehingga $G =
F + c$ (\cref{cor:b1:derivative:monotone}), dan $G(b) - G(a) = F(b)
- F(a) = \int_a^b f$.
\end{proof}

\begin{example}[Kesetangkupan sebelum perhitungan]\label{ex:b1:integration:parity}
Pada \omterm{prop:b1:reals:intervals}{selang} yang setangkup, keparitasannya yang bekerja: jika $f$ ganjil,
maka penyulihan $t \mapsto -t$ mengirim $\int_{-a}^{0} f$ ke
$-\int_0^a f$, sehingga
\[
\int_{-a}^{a} f(t)\,\dd t = 0 ;
\qquad\text{sedangkan jika $f$ genap,}\quad
\int_{-a}^{a} f = 2\int_0^a f .
\]
Jadi $\int_{-1}^{1} \frac{t^3\cos t}{1 + t^4}\,\dd t = 0$ tanpa
primitif yang terlihat (karena integrannya ganjil), dan
$\int_{-\pi}^{\pi} t^2\cos t\,\dd t = 2\int_0^\pi t^2\cos t\,
\dd t$. Jadi periksalah kesetangkupannya sebelum meraih tekniknya: karena
\omterm{thm:b1:integration:def}{integral} tercepat adalah yang tak pernah dihitung.
\end{example}

\begin{example}[Mengenali sebuah turunan pada pandangan pertama]\label{ex:b1:integration:halfangle}
Hitunglah $\displaystyle\int_0^{\pi/2} \frac{\dd x}{1 + \cos x}$.
Kesamaan sudut paruhnya $1 + \cos x = 2\cos^2\frac x2$ mengubah
integrannya menjadi $\frac{1}{2}\bigl(1 + \tan^2\frac
x2\bigr)$, yang persis merupakan \omterm{def:b1:derivative:def}{turunan} $\tan\frac x2$:
\[
\int_0^{\pi/2} \frac{\dd x}{1 + \cos x}
= \Bigl[\tan\frac x2\Bigr]_0^{\pi/2} = \tan\frac\pi4 = 1 .
\]
Jadi tak ada mesin penyulihan yang diperlukan --- hanya refleks
membaca sebuah integran sebagai \omterm{def:b1:derivative:def}{turunan} seseorang, dan teorema fundamentalnya
mengerjakan sisanya. (Adapun perkakas sistematis di balik \omterm{thm:b1:integration:def}{integral}
trigonometri semacam itu, yaitu penyulihan $t = \tan\frac x2$,
termasuk kotak perkakas baku yang dibangun dari
\cref{thm:b1:integration:parts}~(2).)
\end{example}

\begin{example}[Fungsi yang didefinisikan oleh integral]\label{ex:b1:integration:definedby}
Teorema fundamentalnya memproduksi fungsi. Misalkan
\[
F(x) = \int_0^x \eu^{-t^2}\,\dd t .
\]
Tak ada kombinasi fungsi klasik yang berturunan
$\eu^{-t^2}$ (yaitu teorema Liouville, yang diterima); namun $F$ ada,
berkelas $C^1$ dengan $F'(x) = \eu^{-x^2} > 0$, naik tegas, ganjil
(lewat penyulihan $t \mapsto -t$), dan terbatas: karena untuk $x \geq 1$,
\[
F(x) - F(1) = \int_1^x \eu^{-t^2}\dd t
\leq \int_1^x \eu^{-t}\dd t \leq \eu^{-1} ,
\]
sehingga $F \leq F(1) + \eu^{-1} \leq 1 + \eu^{-1}$. (Adapun limitnya yang persis,
$\frac{\sqrt\pi}{2}$, dihitung dengan \omterm{thm:b1:integration:def}{integral} rangkap pada jilid
Tahun ke-3.) Aturan rantai bagi batas yang bergerak: $\frac{\dd}{\dd
x}\int_x^{x^2} \eu^{-t^2}\dd t = 2x\,\eu^{-x^4} - \eu^{-x^2}$.
Inti gagasan penutupnya: bahwa pengintegralan \emph{menciptakan} fungsi baru
dari yang lama, dengan segala sifatnya terbaca dari
integrannya --- jadi primitif yang tak dapat kamu tuliskan tetap
fungsi yang kamu kendalikan sepenuhnya.
\end{example}

\begin{example}[Menaksir tanpa menilai]\label{ex:b1:integration:estimate}
\omterm{thm:b1:integration:def}{Integral} $R_n = \int_0^1 \frac{t^n}{1 + t}\,\dd t$ tak berbentuk
\omterm{def:b1:topology:closed}{tertutup} yang menyenangkan, namun kemonotonannya memakukannya dengan tepat: karena pada
$\intcc{0}{1}$, $\frac12 \leq \frac{1}{1+t} \leq 1$, sehingga
\[
\frac{1}{2(n+1)} = \frac12\int_0^1 t^n\,\dd t
\;\leq\; R_n \;\leq\; \int_0^1 t^n \,\dd t = \frac{1}{n+1} :
\]
jadi orde peluruhannya yang persis ($R_n \sim$ kelipatan $\frac1n$,
dan sesungguhnya $R_n \sim \frac{1}{2n}$) dengan dua baris dan tanpa
antiturunan. Adapun soal akhir pekan bab ini dan bab
berikutnya berjalan persis di atas apitan semacam itu --- jadi naluri pertama
seorang analis di hadapan sebuah \omterm{thm:b1:integration:def}{integral} seharusnya \emph{membatasinya}, lalu
barulah, bila perlu, menghitungnya.
\end{example}

\begin{example}[Nilai rata-rata]\label{ex:b1:integration:average}
\emph{Rata-rata} sebuah $f$ yang \omterm{def:b1:continuity:continuous}{kontinu} pada $\intcc{a}{b}$ adalah
$\frac{1}{b-a}\int_a^b f$. Untuk lengkung sinusnya:
\[
\frac{1}{\pi}\int_0^\pi \sin t\,\dd t
= \frac{1}{\pi}\bigl[-\cos t\bigr]_0^\pi = \frac{2}{\pi}
\approx 0.637 :
\]
jadi satu lengkung positif yang penuh berata-rata bukan $\frac12$ melainkan
$\frac2\pi$ --- karena kurvanya menghabiskan waktu lebih lama di atas daripada segitiga
yang sebanding. Menurut \cref{exo:b1:integration:11} (yaitu teorema nilai rata-rata bagi
\omterm{thm:b1:integration:def}{integral}, dengan $g = 1$), rata-ratanya merupakan sebuah \emph{nilai}: $\sin c =
\frac2\pi$ untuk suatu $c \in \intoo{0}{\pi}$. Dan menurut \omterm{thm:b1:integration:riemann}{jumlah
Riemann} pada bab ini, rata-ratanya merupakan limit rata-rata biasa
atas $n$ cuplikan --- jadi jembatan antara rata-rata diskret
sebuah data dan rata-rata \omterm{def:b1:continuity:continuous}{kontinu} sebuah isyarat, dan begitulah
\omterm{thm:b1:integration:def}{integralnya} masuk ke fisika.
\end{example}

\begin{theorem}[Pengintegralan parsial; penyulihan]\label{thm:b1:integration:parts}
\begin{enumerate}
  \item Jika $u, v$ berkelas $C^1$ pada $\intcc{a}{b}$:\index{pengintegralan
        parsial}
        \[
        \int_a^b u'v = \bigl[uv\bigr]_a^b - \int_a^b uv' .
        \]
  \item Jika $\varphi$ berkelas $C^1$ pada $\intcc{\alpha}{\beta}$ dan $f$
        \omterm{def:b1:continuity:continuous}{kontinu} pada $\varphi(\intcc{\alpha}{\beta})$:
        \index{penyulihan}
        \[
        \int_{\alpha}^{\beta} f\bigl(\varphi(t)\bigr)\,\varphi'(t)\,
        \dd t = \int_{\varphi(\alpha)}^{\varphi(\beta)} f(x)\, \dd x .
        \]
\end{enumerate}
\end{theorem}

\begin{proof}
(1) Karena $(uv)' = u'v + uv'$; integralkanlah atas $\intcc{a}{b}$ lalu terapkan
teorema fundamentalnya pada fungsi $C^1$ yaitu $uv$.

(2) Misalkan $F$ sebuah primitif bagi $f$ pada \omterm{prop:b1:reals:intervals}{selang} petanya
(\cref{thm:b1:integration:ftc}). Maka $(F \circ \varphi)' = (f
\circ \varphi)\,\varphi'$ (menurut aturan rantainya), sehingga kedua ruasnya sama dengan
$F(\varphi(\beta)) - F(\varphi(\alpha))$.
\end{proof}

\begin{example}\label{ex:b1:integration:parts}
$\int_0^1 t\,\eu^t \dd t = \bigl[t\,\eu^t\bigr]_0^1 - \int_0^1
\eu^t\dd t = \eu - (\eu - 1) = 1$. Lalu dengan penyulihan $x =
\sin t$ ($t \in \intcc{0}{\frac\pi2}$):
\[
\int_0^1 \sqrt{1 - x^2}\, \dd x
= \int_0^{\pi/2} \cos t \cdot \cos t \, \dd t
= \int_0^{\pi/2} \frac{1 + \cos 2t}{2}\, \dd t = \frac\pi4 ,
\]
--- yaitu seperempat cakram satuan, sebagaimana dituntut geometrinya. (Jadi pelinearan
dari \cref{met:b1:complex:trig} sedang bekerja.)
\end{example}

\begin{remark}[Jebakan yang lazim dalam kalkulus integral]
(i) \emph{Penyulihannya harus $C^1$ pada seluruh \omterm{prop:b1:reals:intervals}{selangnya}}:
karena penggantian $x = \frac1t$ tak sah melintasi $0$; dan bila diterapkan secara buta
pada $\int_{-1}^{1}\frac{\dd x}{1 + x^2}$ ia ``membuktikan'' bahwa
\omterm{thm:b1:integration:def}{integralnya} sama dengan negatifnya sendiri. Jadi ketika sebuah penyulihan mempunyai
kesingularan, potonglah \omterm{prop:b1:reals:intervals}{selangnya} lebih dulu (dengan Chasles), lalu sulihkan pada
setiap kepingnya, dan barulah gabungkan kembali. (ii) \emph{Primitif
logaritma menuntut nilai mutlak}: karena $\int \frac{\dd x}{x - 2} =
\ln\abs{x - 2} + C$ pada masing-masing sisi $2$ secara tersendiri --- sehingga menulis
$\ln(x - 2)$ pada $\intoo{0}{1}$ berarti menulis logaritma bilangan
yang negatif; dan konstantanya $C$ boleh berbeda pada kedua
sisi kesingularannya. (iii) \emph{\omterm{thm:b1:integration:def}{Integral} yang lenyap tak
membunuh fungsinya}: karena $\int_0^{2\pi}\sin = 0$; jadi kepositifan
integrannya dituntut sebelum menyimpulkan $f = 0$
(\cref{ex:b1:integration:chasleswork}). (iv) \emph{\omterm{thm:b1:integration:riemann}{Jumlah Riemann}
harus terkalibrasi}: karena pada $\frac{b-a}{n}\sum f\bigl(a +
k\frac{b-a}{n}\bigr)$, langkah di luarnya dan titik di dalamnya
harus cocok dengan \omterm{thm:b1:logic:partition}{partisi} yang sama --- adapun galat yang sering terjadi adalah jumlah
$\sum_{k=1}^{n} f\bigl(\frac kn\bigr)$ \emph{tanpa} faktor
$\frac1n$, yang divergen alih-alih konvergen ke $\int_0^1
f$. Jadi daftar periksa sebelum memanggil \cref{thm:b1:integration:riemann}:
keluarkanlah faktor $\frac1n$, tulislah ulang sukunya sebagai $f$ dari $\frac kn$,
namailah $f$ lalu periksalah \omterm{def:b1:continuity:continuous}{kekontinuannya}.
\end{remark}

\begin{example}[Tebak, turunkan, sesuaikan]\label{ex:b1:integration:guessdiff}
Berapakah $\int_1^x (\ln t)^2\,\dd t$? Tebaklah sebuah primitif
berbentuk $t\,P(\ln t)$ dengan $P$ \omterm{def:b1:poly:def}{polinomial} lalu turunkanlah:
\[
\bigl(t\,P(\ln t)\bigr)' = P(\ln t) + P'(\ln t) .
\]
Kita memerlukan $P(u) + P'(u) = u^2$: ambillah $P(u) = u^2 - 2u + 2$
(dengan mencocokkan koefisiennya turun dari $u^2$). Sehingga
\[
\int_1^x (\ln t)^2\,\dd t
= \bigl[t\bigl((\ln t)^2 - 2\ln t + 2\bigr)\bigr]_1^x
= x(\ln x)^2 - 2x\ln x + 2x - 2 ,
\]
yaitu hasil yang selainnya dicapai lewat dua kali \omterm{thm:b1:integration:parts}{pengintegralan parsial}. Inti
gagasan penutupnya: bahwa untuk integran berbentuk
(\omterm{def:b1:poly:def}{polinomial} dalam $\ln t$) atau (\omterm{def:b1:poly:def}{polinomial} dikali $\eu^{\lambda t}$),
primitifnya berbentuk sama --- sehingga menurunkan sebuah tebakan
yang terbentuk mengubah pengintegralan menjadi aljabar linear atas koefisiennya,
yang lebih cepat dan lebih kecil kemungkinan galatnya daripada parsial yang berulang.
\end{example}

\begin{example}[Integral bumerang]\label{ex:b1:integration:boomerang}
Hitunglah $I = \int_0^{\pi/2} \eu^x \cos x\,\dd x$. Integralkanlah secara
parsial dua kali, dengan menurunkan faktor trigonometrinya setiap kalinya:
\[
I = \bigl[\eu^x \sin x\bigr]_0^{\pi/2}
- \int_0^{\pi/2} \eu^x \sin x\,\dd x
= \eu^{\pi/2} - J,
\qquad
J = \bigl[-\eu^x\cos x\bigr]_0^{\pi/2}
+ \int_0^{\pi/2} \eu^x\cos x\,\dd x = 1 + I .
\]
\omterm{thm:b1:integration:def}{Integralnya} sudah kembali ke dirinya sendiri: $I = \eu^{\pi/2} - 1 - I$,
sehingga
\[
I = \frac{\eu^{\pi/2} - 1}{2} .
\]
Inti gagasan penutupnya: bahwa ketika integrannya hasil kali dua
fungsi yang memproduksi dirinya sendiri di bawah pendiferensialan
($\eu^{ax}$, $\cos bx$, $\sin bx$), maka dua kali \omterm{thm:b1:integration:parts}{pengintegralan parsial}
menghasilkan persamaan linear bagi \omterm{thm:b1:integration:def}{integral} yang tak diketahuinya --- jadi pecahkanlah ia
alih-alih mengintegralkannya; dan setara dengan itu, lewatilah
$\eu^{(a + \iu b)x}$ (\cref{ch:b1:complex}) lalu ambillah bagian
realnya. Kedua jalannya memberikan jawaban yang sama, dan memeriksa bahwa demikianlah
halnya merupakan uji kewarasan yang cuma-cuma.
\end{example}

\section{Jumlah Riemann}

\begin{theorem}[Jumlah Riemann]\label{thm:b1:integration:riemann}
Misalkan $f$ \omterm{def:b1:continuity:continuous}{kontinu} pada $\intcc{a}{b}$. Maka\index{jumlah Riemann}
\[
S_n = \frac{b - a}{n} \sum_{k=0}^{n-1}
f\Bigl(a + k\,\frac{b-a}{n}\Bigr)
\xrightarrow[n \to \infty]{} \int_a^b f(t)\, \dd t ,
\]
dan demikian pula dengan sebarang titik penilaian di dalam subselangnya.
\end{theorem}

\begin{proof}
Di sini $S_n$ merupakan \omterm{thm:b1:integration:def}{integral} \omterm{def:b1:integration:step}{fungsi tangga} $\varphi_n$ yang sama dengan
$f(a + k\frac{b-a}{n})$ pada subselang ke-$k$. Diberikan
$\varepsilon > 0$, \omterm{def:b1:continuity:uniform}{kekontinuan seragamnya} (menurut Heine) menyediakan $\delta$; lalu untuk
$n > \frac{b-a}{\delta}$, setiap titik sebuah subselang berjarak
$\delta$ dari titik penilaiannya, sehingga $\abs{f - \varphi_n} \leq
\varepsilon$ pada $\intcc{a}{b}$, jadi
\[
\Bigl| \int_a^b f - S_n \Bigr|
= \Bigl| \int_a^b (f - \varphi_n) \Bigr|
\leq (b-a)\,\varepsilon . \qedhere
\]
\end{proof}

\begin{omfigure}
\begin{tikzpicture}
\begin{axis}[omaxis, width=8.6cm, height=5.4cm,
    xmin=0, xmax=1.15, ymin=0, ymax=1.7,
    xtick={0,1}, ytick=\empty]
  \draw[fill=omDef!15, draw=omDef!60, thin]
    (axis cs:0,0)      rectangle (axis cs:0.125,0.4)
    (axis cs:0.125,0)  rectangle (axis cs:0.25,0.4156)
    (axis cs:0.25,0)   rectangle (axis cs:0.375,0.4625)
    (axis cs:0.375,0)  rectangle (axis cs:0.5,0.5406)
    (axis cs:0.5,0)    rectangle (axis cs:0.625,0.65)
    (axis cs:0.625,0)  rectangle (axis cs:0.75,0.7906)
    (axis cs:0.75,0)   rectangle (axis cs:0.875,0.9625)
    (axis cs:0.875,0)  rectangle (axis cs:1,1.1656);
  \addplot[omThm, very thick, domain=0:1.05, samples=60]
    {0.4 + x^2};
\end{axis}
\end{tikzpicture}

{\small Sebuah \omterm{thm:b1:integration:riemann}{jumlah Riemann} kiri dengan $n = 8$ persegi panjang: bahwa ketika jalanya
menyusut, \omterm{def:b1:continuity:uniform}{kekontinuan seragamnya} memaksa luas tangganya menuju
$\int_a^b f$.}
\end{omfigure}

\begin{example}\label{ex:b1:integration:riemannexample}
$\displaystyle\sum_{k=1}^{n} \frac{1}{n + k} = \frac 1n
\sum_{k=1}^{n} \frac{1}{1 + k/n} \xrightarrow[n\to\infty]{}
\int_0^1 \frac{\dd x}{1 + x} = \ln 2$: yaitu limit yang tak kasatmata bagi
batas elementer, tetapi bening sebagai \omterm{thm:b1:integration:riemann}{jumlah Riemann}.
\end{example}

\begin{example}[Jumlah Riemann kedua, beserta kalibrasinya]\label{ex:b1:integration:riemann2}
Carilah $\lim_{n\to\infty} \sum_{k=1}^{n} \dfrac{n}{(n+k)^2}$.
Kalibrasikanlah:
\[
\sum_{k=1}^{n} \frac{n}{(n+k)^2}
= \frac{1}{n}\sum_{k=1}^{n} \frac{n^2}{(n+k)^2}
= \frac1n \sum_{k=1}^{n} \frac{1}{\bigl(1 + \frac kn\bigr)^2} ,
\]
yaitu \omterm{thm:b1:integration:riemann}{jumlah Riemann} bagi $f(x) = \frac{1}{(1+x)^2}$ yang \omterm{def:b1:continuity:continuous}{kontinu} pada
$\intcc{0}{1}$: sehingga limitnya adalah
\[
\int_0^1 \frac{\dd x}{(1 + x)^2}
= \Bigl[-\frac{1}{1+x}\Bigr]_0^1 = \frac12 .
\]
Inti gagasan penutupnya: bahwa seluruh seninya ada pada baris tengahnya ---
paksalah sukunya menjadi bentuk $f(\frac kn)$ dengan ongkos
mengeluarkan tepat satu faktor $\frac1n$; dan begitu bentuknya
tepat, teoremanya mengerjakan analisisnya sedangkan teorema fundamentalnya
mengerjakan aritmetikanya.
\end{example}

\begin{remark}[Di mana integralnya bekerja berikutnya]
Konstruksi bab ini masing-masing mempunyai kelanjutan. Adapun \omterm{thm:b1:integration:riemann}{jumlah Riemann}
kembali pada \cref{ch:b1:series} sebagai jembatan antara deret dan
\omterm{thm:b1:integration:def}{integral} (yaitu perbandingan $\sum \frac{1}{n^\alpha}$ dengan $\int
\frac{\dd t}{t^\alpha}$); sedangkan sisa \omterm{thm:b1:integration:def}{integralnya} merupakan bentuk paling
tajam rumus Taylor pada \cref{ch:b1:taylor}; dan
definisi berbasis $\sup$-nya merupakan purwarupa bagi \omterm{thm:b1:integration:def}{integral}
Lebesgue pada jilid Tahun ke-3, yang di situ ketiga sifat yang sama
(kelinearan, kemonotonan, dan sebuah teorema kekonvergenan) dibangun kembali pada
kelas fungsi yang jauh lebih besar. Adapun soal akhir pekan di bawah
mengubah \omterm{thm:b1:integration:parts}{pengintegralan parsial} menjadi aritmetika: yaitu keirasionalan
$\pi^2$.
\end{remark}

\section{Latihan}

\begin{exercise}[$\star$]\label{exo:b1:integration:1}
Hitunglah: $\displaystyle\int_0^1 \frac{\dd x}{x^2 - 4}$
\emph{(lewat pecahan parsial, \cref{ch:b1:fractions})};
$\displaystyle\int_1^{\eu} \ln t \,\dd t$;
$\displaystyle\int_0^{\pi} t \sin t\, \dd t$.
\end{exercise}

\begin{exercise}[$\star$]\label{exo:b1:integration:2}
Hitunglah $\displaystyle\int_0^{1} \frac{t}{(t^2+1)^2}\,\dd t$
(lewat penyulihan) dan
$\displaystyle\int_0^{\pi/2} \cos^3 t\, \dd t$ (dengan menulis $\cos^3 =
\cos(1 - \sin^2)$).
\end{exercise}

\begin{exercise}[$\star$]\label{exo:b1:integration:3}
Carilah limitnya, sebagai \omterm{thm:b1:integration:riemann}{jumlah Riemann}:
\[
u_n = \sum_{k=1}^{n} \frac{n}{n^2 + k^2},
\qquad
v_n = \frac{1}{n}\sqrt[n]{\frac{(2n)!}{n!\,n^n}}
\ \emph{(ambillah logaritmanya)}.
\]
\end{exercise}

\begin{exercise}[$\star$]\label{exo:b1:integration:4}
Misalkan $f$ \omterm{def:b1:continuity:continuous}{kontinu} pada $\intcc{0}{1}$. Hitunglah $\lim_{n\to\infty}
\int_0^1 x^n f(x)\,\dd x$. \emph{(Potonglah $\intcc{0}{1}$ di $1 -
\delta$.)}
\end{exercise}

\begin{exercise}[$\star\star$]\label{exo:b1:integration:5}
(Cauchy--Schwarz) Untuk $f, g$ yang \omterm{def:b1:continuity:continuous}{kontinu} pada $\intcc{a}{b}$, buktikanlah
\[
\Bigl(\int_a^b fg\Bigr)^{\!2} \leq \int_a^b f^2 \cdot \int_a^b g^2 ,
\]
dengan menjabarkan $\int_a^b (f + \lambda g)^2 \geq 0$ sebagai kuadratik dalam
$\lambda$. Kapankah ia menjadi kesamaan?
\end{exercise}

\begin{exercise}[$\star\star$]\label{exo:b1:integration:6}
Misalkan $f$ \omterm{def:b1:continuity:continuous}{kontinu} pada $\R$, dan periodik-$T$. Buktikan bahwa $\int_a^{a+T}
f$ tak bergantung pada $a$, dan bahwa $\frac1x \int_0^x f(t)\,\dd t
\to \frac 1T \int_0^T f$ ketika $x \to +\infty$.
\end{exercise}

\begin{exercise}[$\star\star$]\label{exo:b1:integration:7}
Untuk $f$ yang \omterm{def:b1:continuity:continuous}{kontinu} pada $\intcc{0}{1}$ dengan $\int_0^1 f = \frac12$,
buktikan bahwa $f$ mempunyai titik tetap di $\intcc{0}{1}$. \emph{(Integralkanlah
$f(x) - x$ lalu pakailah kepositifan tegasnya,
\cref{thm:b1:integration:props}~(4), lewat kontraposisinya
yang digabung dengan teorema nilai antara.)}
\end{exercise}

\begin{exercise}[$\star\star\star$]\label{exo:b1:integration:8}
(\omterm{thm:b1:integration:def}{Integral} Wallis) Misalkan $W_n = \int_0^{\pi/2} \sin^n t\,\dd t$.
\begin{enumerate}
  \item Buktikan rekurensi $n W_n = (n-1) W_{n-2}$ ($n \geq 2$) lewat
        parsial, lalu hitunglah $W_0, W_1$, kemudian $W_{2p}$ dan $W_{2p+1}$
        dalam bentuk \omterm{def:b1:topology:closed}{tertutup}.
  \item Buktikan bahwa $(W_n)$ turun dengan $\frac{W_{n+1}}{W_n}
        \to 1$; buktikan bahwa besaran $(n+1)\,W_{n+1} W_n$ bersifat
        konstan, yang sama dengan $\frac\pi2$; lalu simpulkanlah kesetaraan
        $W_n \sim \sqrt{\dfrac{\pi}{2n}}$.
\end{enumerate}
\end{exercise}

\begin{exercise}[$\star\star\star$]\label{exo:b1:integration:9}
(Niven: $\pi$ irasional) Andaikan $\pi = \frac ab$ dengan $a, b \in
\N^*$, lalu tetapkan, untuk $n$ yang akan dipilih,
\[
P(x) = \frac{x^n (a - bx)^n}{n!},
\qquad
I_n = \int_0^{\pi} P(x)\sin x\, \dd x .
\]
\begin{enumerate}
  \item Buktikan bahwa $0 < I_n \leq \pi\,\frac{(\pi a)^n}{n!}$, yang
        bernilai $< 1$ untuk $n$ yang besar.
  \item Buktikan bahwa $P$ beserta semua \omterm{def:b1:derivative:def}{turunannya} mengambil nilai
        \emph{bulat} di $0$ dan di $\pi = \frac ab$. \emph{(Lewat penjabaran
        binomialnya: karena koefisien $P$ dikali $k!$ bersifat bulat
        untuk $k \geq n$; dan $P(\pi - x) = P(x)$.)}
  \item Tetapkanlah $Q = P - P'' + P^{(4)} - \dots$ (yaitu jumlah yang hingga). Periksalah
        bahwa $\bigl(Q'\sin x - Q\cos x\bigr)' = P \sin x$, lalu simpulkanlah
        bahwa $I_n = Q(\pi) + Q(0)$ bersifat bulat.
  \item Tutuplah kesimpulannya.
\end{enumerate}
\end{exercise}

\begin{exercise}[$\star\star\star$]\label{exo:b1:integration:10}
Misalkan $f$ berkelas $C^1$ pada $\intcc{a}{b}$. Buktikan limit bertipe
Riemann--Lebesgue
\[
\int_a^b f(t)\sin(\lambda t)\,\dd t \xrightarrow[\lambda \to
+\infty]{} 0
\]
lewat \omterm{thm:b1:integration:parts}{pengintegralan parsial}. Lalu buktikanlah ia lagi untuk $f$ yang sekadar
\omterm{def:b1:continuity:continuous}{kontinu}, lewat hampiran seragam dengan \omterm{def:b1:integration:step}{fungsi tangga}
(\cref{thm:b1:integration:approx}).
\end{exercise}

\begin{exercise}[$\star\star$]\label{exo:b1:integration:11}
(Teorema nilai rata-rata bagi \omterm{thm:b1:integration:def}{integral}) Misalkan $f, g$ \omterm{def:b1:continuity:continuous}{kontinu} pada
$\intcc{a}{b}$ dengan $g \geq 0$. Buktikan bahwa ada $c \in
\intcc{a}{b}$ dengan
\[
\int_a^b f(t)\,g(t)\,\dd t = f(c)\int_a^b g(t)\,\dd t ,
\]
lalu tunjukkanlah lewat sebuah contoh bahwa hipotesis $g \geq 0$ tak dapat
dilepaskan.
\end{exercise}

\begin{exercise}[$\star\star\star$]\label{exo:b1:integration:12}
(Momen memaksa nol) Misalkan $f$ \omterm{def:b1:continuity:continuous}{kontinu} pada $\intcc{a}{b}$
dengan
\[
\int_a^b f(t)\,t^k\,\dd t = 0 \qquad \text{untuk } k = 0, 1,
\dots, n .
\]
Buktikan bahwa $f$ lenyap di $n + 1$ titik berbeda pada
$\intoo{a}{b}$. \emph{(Jika $f$ berubah tanda hanya di $z_1 < \dots
< z_m$ dengan $m \leq n$, integralkanlah $f$ terhadap $P(t) = (t - z_1)
\cdots (t - z_m)$ lalu pakailah kepositifan tegasnya.)}
\end{exercise}

\begin{remark}[Cakrawala di dalam jilid ini]
Tiga bab ke depan bersandar langsung pada bab ini.
\Cref{ch:b1:taylor} membawa sisa \omterm{thm:b1:integration:def}{integralnya} --- karena yang
paling tajam di antara ketiga rumus Taylornya adalah \omterm{thm:b1:integration:parts}{pengintegralan parsial}
yang diulang $n$ kali. \Cref{ch:b1:series} mengubah apitan
jumlah oleh \omterm{thm:b1:integration:def}{integral} menjadi uji yang menentukan bagi $\sum n^{-\alpha}$,
dan soal akhir pekannya menghaluskan apitan itu menjadi konstanta
Euler. \Cref{ch:b1:curves} membuat \omterm{thm:b1:integration:def}{integralnya} geometris: karena
panjang sebuah busur berparameter adalah $\int \sqrt{x'(t)^2 +
y'(t)^2}\,\dd t$, yaitu \omterm{thm:b1:integration:def}{integral} fungsi yang \omterm{def:b1:continuity:continuous}{kontinu} pada sebuah
ruas --- persis objek yang dibangun di sini, tanpa memerlukan teori
takwajar apa pun. Adapun fakta yang paling banyak dipakai ulang justru yang paling rendah hati:
$\bigl|\int f\bigr| \leq (b - a)\sup\abs{f}$, yaitu ketaksamaan
yang mengubah setiap taksiran titik demi titik menjadi taksiran \omterm{thm:b1:integration:def}{integral}.
\end{remark}

\section{Soal: Mesin keirasionalan lewat integral}

\begin{problem}[{Soal akhir pekan --- $\eu$ dan $\pi^2$ bersifat
irasional, $\eu$ sampai enam desimal, dan $\frac{22}{7} > \pi$ beserta
buktinya}]
\label{pb:b1:integration:1}
Satu mekanisme menggerakkan seluruh soal ini: \emph{bahwa ungkapan yang
harus menjadi bilangan bulat positif, namun terbukti lebih kecil daripada $1$,
tak mungkin ada}. \Cref{exo:b1:integration:9} (Niven) menjalankannya sekali untuk
membuktikan $\pi \notin \Q$; sedangkan di sini kita mengindustrikannya. Mesinnya
memerlukan tiga bagian: masukan \emph{kebulatan} (yaitu nilai titik ujung
\omterm{def:b1:poly:def}{polinomial} yang terpilih baik), masukan \emph{kekecilan} (yaitu faktor
$\frac{1}{n!}$ yang menggilas \omterm{thm:b1:integration:def}{integralnya}), dan sebuah \emph{jembatan}
(yaitu \omterm{thm:b1:integration:parts}{pengintegralan parsial}) yang menghubungkan keduanya. Kita membuktikan bahwa $\eu$ bersifat
irasional lalu menghitungnya dengan galat yang bersertifikat, membuktikan teorema
Legendre yang lebih tajam bahwa $\pi^2$ irasional, lalu berakhir dengan
\omterm{thm:b1:integration:def}{integral} paling memikat dalam analisis: $\int_0^1 \frac{x^4(1 -
x)^4}{1 + x^2}\dd x = \frac{22}{7} - \pi$, yang mengapit $\pi$
dengan tangan.\index{keirasionalan!pi kuadrat@$\pi^2$}\index{keirasionalan!e@$\eu$}

\medskip
\noindent\textbf{Bagian I --- Bahan bakarnya.}
\begin{enumerate}
  \item Buktikan bahwa $\dfrac{c^{\,n}}{n!} \to 0$ untuk setiap $c
        > 0$ yang tetap \emph{(karena di luar $n \geq 2c$, setiap langkahnya sekurang-kurangnya
        memaruhkan sukunya)}.
  \item (\omterm{thm:b1:integration:def}{Integral} Beta) Buktikan, lewat induksi pada $m$ dengan
        \omterm{thm:b1:integration:parts}{pengintegralan parsial}:
        \[
        \int_0^1 x^{\,k}\,(1 - x)^{\,m}\,\dd x =
        \frac{k!\,m!}{(k + m + 1)!}
        \qquad (k, m \in \N).
        \]
  \item Simpulkan $\displaystyle\int_0^1 \bigl(x(1-x)\bigr)^n \dd x
        = \frac{1}{(2n+1)\binom{2n}{n}}$, dan --- dengan membandingkannya dengan
        batas $x(1 - x) \leq \frac14$ --- taksiran
        $\binom{2n}{n} \geq \dfrac{4^n}{2n+1}$, yang cocok dengan
        $\binom{2n}{n}^{1/n} \to 4$ dari \cref{pb:b1:seq:1}.
  \item Buktikan lema kekecilan yang dipakai dua kali di bawah: bahwa untuk setiap
        $g > 0$ yang \omterm{def:b1:continuity:continuous}{kontinu} pada $\intoo{0}{1}$,
        \[
        0 < \int_0^1 \bigl(x(1-x)\bigr)^n g(x)\,\dd x
        \leq \frac{\sup_{\intcc{0}{1}}\abs g}{4^{\,n}} .
        \]
\end{enumerate}

\medskip
\noindent\textbf{Bagian II --- $\eu$: keirasionalannya, lalu enam
desimalnya.} Tetapkanlah $A_n = \displaystyle\int_0^1 x^n \eu^x \dd x$.
\begin{enumerate}[resume]
  \item Hitunglah $A_0$ dan $A_1$, buktikan rekurensi $A_n = \eu
        - n A_{n-1}$, beserta batasnya $0 < A_n \leq
        \dfrac{\eu}{n+1}$.
  \item Tunjukkan lewat induksi bahwa $A_n = \alpha_n + \beta_n \eu$
        dengan $\alpha_n, \beta_n \in \Z$.
  \item Simpulkan bahwa $\eu$ irasional \emph{(karena jika $\eu = \frac
        pq$, maka $q A_n$ merupakan bilangan bulat yang terperangkap di
        $\intoo{0}{1}$ untuk $n$ yang besar)}. Bandingkanlah dengan bukti
        \cref{exo:b1:seq:9}: yaitu \omterm{def:b1:topology:closure}{penutup} yang sama, dengan bahan bakar yang berbeda.
  \item Buktikan, lewat induksi dan \omterm{thm:b1:integration:parts}{pengintegralan parsial}, rumus
        sisanya yang persis
        \[
        \eu = \sum_{k=0}^{n} \frac{1}{k!} + R_n,
        \qquad
        R_n = \frac{1}{n!}\int_0^1 (1 - t)^n\,\eu^{\,t}\,\dd t,
        \qquad
        \frac{1}{(n+1)!} \leq R_n \leq \frac{\eu}{(n+1)!} .
        \]
  \item Ambillah $n = 9$: batasilah $R_9$ dengan memakai $\eu < 2.75$ (dari
        $b_2 = 2.75$ pada \cref{ex:b1:seq:e}), nilailah jumlahnya,
        lalu simpulkan apitan bersertifikatnya $2.7182818 \leq
        \eu \leq 2.7182823$ --- jadi enam desimal, $\eu \approx
        2.718282$, beserta buktinya.
\end{enumerate}

\medskip
\noindent\textbf{Bagian III --- Teorema Legendre: bahwa $\pi^2$
irasional.} Misalkan $f(x) = \dfrac{x^n (1 - x)^n}{n!}$, dan andaikan
$\pi^2 = \frac ab$ dengan $a, b \in \N^*$.
\begin{enumerate}[resume]
  \item Tunjukkan $f(1 - x) = f(x)$ dan $0 < f \leq
        \dfrac{1}{4^n\,n!}$ pada $\intoo{0}{1}$.
  \item Tunjukkan bahwa $f^{(k)}(0)$ dan $f^{(k)}(1)$ bersifat bulat untuk
        setiap $k \geq 0$ \emph{(jabarkanlah $x^n(1-x)^n$ dengan koefisien
        bulat; lalu $\frac{k!}{n!} \in \Z$ untuk $k \geq n$;
        kemudian pakailah kesetangkupannya)}.
  \item Definisikanlah
        \[
        G = b^{\,n} \sum_{k=0}^{n} (-1)^k\,
        \pi^{2n - 2k} f^{(2k)} .
        \]
        Tunjukkan bahwa $G(0)$ dan $G(1)$ bersifat bulat \emph{(karena masing-masing
        $b^n \pi^{2n-2k} = a^{\,n-k}\,b^{\,k}$)}.
  \item Periksalah teleskopnya $G'' + \pi^2 G = b^n \pi^{2n+2} f
        = \pi^2 a^n f$, lalu
        \[
        \frac{\dd}{\dd x}\Bigl(G'(x)\sin \pi x - \pi\,G(x)\cos
        \pi x\Bigr) = \pi^2 a^n f(x)\sin \pi x .
        \]
  \item Integralkanlah atas $\intcc{0}{1}$ lalu simpulkanlah
        \[
        \pi\,a^n \int_0^1 f(x)\sin(\pi x)\,\dd x = G(0) + G(1)
        \in \Z ,
        \]
        yaitu bilangan bulat yang \emph{positif} yang terbatas oleh $\dfrac{\pi
        a^n}{4^n\,n!}$.
  \item Tutuplah dengan pertanyaan 1 bahwa $\pi^2$ bersifat irasional
        (Legendre, 1794), dan bahwa ini memperkuat
        \cref{exo:b1:integration:9}: mengapakah keirasionalan
        $\pi^2$ mengakibatkan keirasionalan $\pi$, dan tidak sebaliknya?
\end{enumerate}

\medskip
\noindent\textbf{Bagian IV --- Memahami mesinnya.}
\begin{enumerate}[resume]
  \item Tunjukkanlah kedua daya yang berlawanan (yaitu kebulatan data titik
        ujungnya; dan kekecilan analitik \omterm{thm:b1:integration:def}{integralnya}) beserta
        jembatannya, pada Bagian II dan III. Lalu jelaskan mengapa faktor
        $\frac{1}{n!}$ pada $f$ menjadi titik gentingnya: bahwa jika ia dibuang,
        kebulatannya bertahan, tetapi ketaksamaan mana yang mati, dan untuk
        pecahan $\frac ab$ yang diklaim yang mana buktinya lalu
        gagal?
  \item Keefektifannya: andaikan seseorang mengklaim $\pi^2 = \frac ab$
        dengan $a \leq 10$. Tunjukkan bahwa pertentangannya sudah
        mendarat di $n = 7$: hitunglah $\pi\,(10/4)^7/7! \approx
        0.38 < 1$. Jadi mesinnya tak sekadar membantah; ia
        membantah pada tahap yang tetap dan terhitung.
  \item Periksalah kewarasan jembatannya tanpa syarat: buktikan lewat dua kali
        \omterm{thm:b1:integration:parts}{pengintegralan parsial} bahwa
        \[
        \int_0^1 x(1 - x)\sin(\pi x)\,\dd x = \frac{4}{\pi^3},
        \]
        lalu rukunkanlah dengan pertanyaan 14 di $n = 1$ (dengan menjaga
        $\pi^2$ tetap simbolis: karena kesamaan teleskopnya berbunyi $\pi^3
        \int_0^1 f_1 \sin \pi x = -(f_1''(0) + f_1''(1)) = 4$).
  \item Apakah yang membuat $\eu^x$ dan $\sin \pi x$ layak menjadi
        inti mesinnya? Kenalilah sifatnya (bahwa masing-masingnya memenuhi
        persamaan diferensial linear berkoefisien
        konstan, sehingga \omterm{thm:b1:integration:parts}{pengintegralan parsial} yang berulang berdaur
        kembali ke awalnya), lalu namailah \omterm{def:b1:topology:closure}{perbatasannya}: bahwa mesin yang
        sama, yang dihaluskan oleh Hermite dan Lindemann, membuktikan bahwa $\eu$
        dan $\pi$ bersifat \emph{transenden} --- yang di luar jilid ini.
\end{enumerate}

\medskip
\noindent\textbf{Bagian V --- $\frac{22}{7}$ lawan $\pi$, dan
moralnya.}
\begin{enumerate}[resume]
  \item Tegakkanlah pembagian \omterm{def:b1:poly:def}{polinomialnya}
        \[
        \frac{x^4(1-x)^4}{1 + x^2}
        = x^6 - 4x^5 + 5x^4 - 4x^2 + 4 - \frac{4}{1 + x^2},
        \]
        lalu simpulkanlah kesamaan yang termasyhur
        \[
        \int_0^1 \frac{x^4 (1-x)^4}{1 + x^2}\,\dd x
        = \frac{22}{7} - \pi .
        \]
  \item Integrannya positif: simpulkanlah $\pi <
        \frac{22}{7}$. Lalu, dengan membatasi $\frac{1}{1+x^2}$ di antara
        $\frac12$ dan $1$ dan memakai $\int_0^1 (x(1-x))^4 =
        \frac{1}{630}$ (pertanyaan 3), buktikanlah
        \[
        \frac{22}{7} - \frac{1}{630} \;\leq\; \pi \;\leq\;
        \frac{22}{7} - \frac{1}{1260} ,
        \]
        yakni $3.14126 \leq \pi \leq 3.14207$: jadi dua desimal
        yang benar, dengan tangan.
  \item Samaratakanlah: dengan membagi $x^{4m}(1-x)^{4m}$ dengan $1 + x^2$,
        tunjukkan bahwa sisanya adalah konstanta $(-4)^m$ \emph{(bekerjalah
        modulo $x^2 + 1$: karena $(1-x)^2 \equiv -2x$)}, lalu simpulkanlah
        adanya bilangan rasional $r_m$ dengan
        \[
        \abs{\pi - r_m} \leq 4^{\,1 - 5m} ,
        \]
        lalu periksalah bahwa $m = 1$ memproduksi ulang pertanyaan 20--21.
  \item Hadapkanlah bilangan rasional ini dengan teori hampiran
        pada \cref{pb:b1:derivative:1}: hitunglah $\abs{\pi -
        \frac{22}{7}} \approx 1.26\cdot10^{-3}$ terhadap
        jaminan Dirichlet $\frac{1}{49}$, lalu kutiplah
        $\abs{\pi - \frac{355}{113}} \approx 2.7\cdot10^{-7}$
        terhadap $\frac{1}{113^2} \approx 7.8\cdot10^{-5}$:
        jadi ada hampiran rasional yang luar biasa baik bagi
        $\pi$ --- yang selaras, karena $\pi$ tak diketahui
        buruk penghampirannya.
  \item (Perangkap bilangan bulat, yang diabstraksikan) Buktikanlah lema yang
        menyatukan segalanya: \emph{bahwa jika $x \in \R$ dan ada
        bilangan bulat $a_n, b_n$ dengan $0 < \abs{a_n + b_n x} \to 0$,
        maka $x$ irasional.} Daftarkanlah kejadiannya pada soal
        ini, pada \cref{exo:b1:integration:9}, pada
        \cref{exo:b1:seq:9}, dan pada \cref{pb:b1:derivative:1}.
  \item Sintesis, satu kalimat untuk masing-masing: (i) ketiga bagian
        mesinnya dan di mana masing-masingnya tinggal pada kotak perkakas bab
        ini; (ii) apakah yang disumbangkan \omterm{thm:b1:integration:def}{integralnya} yang tak dapat
        disumbangkan teorema nilai rata-rata pada \cref{pb:b1:derivative:1};
        (iii) daftar hasil yang terpetik (yaitu dua
        keirasionalan, satu konstanta berdesimal enam, satu
        apitan bagi $\pi$, dan satu batas binomial); (iv) adapun
        \omterm{def:b1:topology:closure}{perbatasannya} (Hermite, Lindemann; beserta perangkap yang sama, yang dijalankan pada
        $\zeta(2)$ dan $\zeta(3)$, dalam aritmetika abad
        kedua puluh).

\end{enumerate}
\end{problem}
